\documentclass[12pt,a4paper]{article}

% --------------------
% Packages
% --------------------
\usepackage{geometry}
\geometry{margin=2.5cm}

\usepackage{amsmath,amssymb}
\usepackage{physics}
\usepackage{graphicx}
\usepackage{setspace}
\usepackage{hyperref}
\usepackage{cite}
\usepackage{indentfirst}

\setlength{\parindent}{2em}
\setlength{\parskip}{0.5em}

% --------------------
% Title Information
% --------------------
\title{\textbf{时间调制石墨烯–金纳米管耦合等离激元系统的设计与研究方案}}
\author{}
\date{}

% --------------------
% Document
% --------------------
\begin{document}

\maketitle
\onehalfspacing

% ==================================================
\section{研究背景与研究意义}

等离激元（plasmonics）体系由于其能够在亚波长尺度内显著增强光–物质相互作用，在光学传感、光调制以及纳米光子器件中具有重要的研究价值。传统等离激元研究主要集中于静态性质，例如局域表面等离激元共振（Localized Surface Plasmon Resonance, LSPR）的共振峰位、线宽及其对周围介电环境的响应。然而，这类静态研究在功能性与信息维度上存在一定局限。

近年来，\textbf{时间调制等离激元系统（temporally modulated plasmonic systems）}逐渐受到关注。通过在时间域中主动调控体系参数（如介电常数、载流子密度或边界条件），可以在频域中引入新的自由度，例如调制带宽、相位延迟、谐波产生以及非线性响应等。这一研究方向不仅拓展了等离激元物理的内涵，也为主动光子器件和动态光学系统提供了新的设计范式。

石墨烯作为一种二维材料，其载流子浓度和费米能级可以通过外加电场实现连续调控，在可见至近红外波段表现出可调的面电导特性。将石墨烯引入金属纳米结构等离激元体系中，使其作为\textbf{主动调制单元}，而金纳米管（或金纳米棒等各向异性金纳米结构）作为\textbf{等离激元谐振腔}，为构建时间调制等离激元系统提供了可行的物理平台。

因此，本项目拟构建并研究一种\textbf{时间调制石墨烯–金纳米管耦合等离激元系统}，重点关注其时间域与频域中的动态光学响应，而非传统意义上的静态 LSPR 传感性能。

% ==================================================
\section{研究目标与核心问题}

\subsection{研究目标}

本项目的总体目标是：

\begin{quote}
通过外加周期性调制手段，对石墨烯–金纳米管耦合体系中的等离激元响应进行动态调控，并对其时间域与频域特征进行系统而定量的实验研究。
\end{quote}

\subsection{核心科学问题}

本项目重点围绕以下三个核心问题展开：

\begin{enumerate}
  \item \textbf{可调性问题}：外部时间调制（如周期性栅压）是否能够引起可重复、可控的光学响应变化，其调制幅度与稳定性如何？
  \item \textbf{动态响应问题}：在不同调制频率下，等离激元响应的幅度和相位特性如何，体系的有效带宽与时间常数是多少？
  \item \textbf{耦合机制问题}：所观测到的调制信号是否确实来源于石墨烯电学性质变化引起的等离激元边界条件改变，而非热效应、离子迁移或表面吸附等非目标效应？
\end{enumerate}

% ==================================================
\section{系统设计与物理原理}

\subsection{系统总体架构}

本项目采用\textbf{电调制石墨烯–金纳米结构等离激元系统}作为主要研究对象，其基本结构包括：

\begin{itemize}
  \item 重掺杂硅基底（作为背栅电极）；
  \item 二氧化硅绝缘层（栅介质）；
  \item 单层或少层石墨烯薄膜；
  \item 超薄介质间隔层（可选）；
  \item 金纳米管（或金纳米棒）等离激元纳米结构。
\end{itemize}

通过对石墨烯施加周期性栅压，实现其费米能级与面电导的时间调制，从而动态改变其与金纳米结构之间的等离激元耦合状态。

\subsection{时间调制机理}

在外加时间变化的栅压 $V_g(t)$ 作用下，石墨烯的载流子浓度 $n(t)$ 随时间变化，其费米能级 $E_F(t)$ 及频率相关面电导 $\sigma(\omega,t)$ 发生调制。作为二维导电层，石墨烯的面电导直接影响电磁场在石墨烯–金纳米结构界面处的边界条件。

由于金纳米管支持局域等离激元共振，其共振特性对局域电磁环境高度敏感，因此石墨烯电导的时间变化将导致：

\begin{itemize}
  \item 等离激元共振强度的周期性变化；
  \item 共振峰位与有效线宽的微小调制；
  \item 固定波长下散射或透射强度的时间调制。
\end{itemize}

在本项目中，等离激元共振主要作为\textbf{信号放大与读出通道}，而非研究对象本身。

% ==================================================
\section{实验方案与器件设计}

\subsection{器件结构与制备思路}

本项目拟采用背栅型石墨烯器件结构，其主要参数如下：

\begin{itemize}
  \item 基底：重掺杂硅；
  \item 栅介质：$90\,\mathrm{nm}$ 或 $285\,\mathrm{nm}$ 厚的 $\mathrm{SiO_2}$；
  \item 石墨烯：CVD 生长并转移至基底；
  \item 金纳米管：通过表面修饰方法固定于石墨烯表面。
\end{itemize}

为调控石墨烯与金纳米结构之间的耦合强度，可在二者之间引入 $1$--$3\,\mathrm{nm}$ 的超薄介质间隔层（如 ALD 生长的 $\mathrm{Al_2O_3}$），并将其厚度作为重要实验变量。

\subsection{调制方式}

本项目采用电调制作为主要调制方式，其形式为：
\begin{equation}
V_g(t) = V_0 + V_m \sin(2\pi f_m t),
\end{equation}
其中 $f_m$ 为调制频率，$V_m$ 为调制幅度。

% ==================================================
\section{光学读出与测量方法}

\subsection{读出策略}

本项目采用\textbf{固定波长强度读出}作为主要测量方式，具体流程如下：

\begin{enumerate}
  \item 测量体系的静态光谱，确定等离激元共振峰；
  \item 选择位于共振谱线斜率最大处的工作波长 $\lambda_{\mathrm{work}}$；
  \item 在固定 $\lambda_{\mathrm{work}}$ 下记录光学强度随时间变化的信号 $I(t)$。
\end{enumerate}

\subsection{数据采集与处理}

实验中将同步记录 $I(t)$ 与 $V_g(t)$，并对 $I(t)$ 进行快速傅里叶变换（FFT），分析其频域特征。通过扫描调制频率，获取体系的幅频与相频响应。

% ==================================================
\section{定量分析指标}

为定量评估时间调制效果，定义以下关键指标：

\begin{enumerate}
  \item \textbf{调制深度}
  \begin{equation}
  M = \frac{\Delta I}{I_0}.
  \end{equation}

  \item \textbf{等效共振调制幅度}
  \begin{equation}
  \Delta \lambda_{\mathrm{eq}} \approx 
  \frac{\Delta I}{\left| \dv{I}{\lambda} \right|_{\lambda_{\mathrm{work}}}}.
  \end{equation}

  \item \textbf{频率响应与时间常数}：通过幅频与相频曲线拟合得到体系的截止频率 $f_c$ 及时间常数 $\tau$。

  \item \textbf{非线性与谐波响应}：通过频谱中高阶谐波分量评估体系的非线性行为。
\end{enumerate}

% ==================================================
\section{对照实验与误差控制}

为排除非目标效应，本项目将设计以下对照实验：

\begin{itemize}
  \item 无金纳米结构的石墨烯器件；
  \item 无石墨烯的金纳米结构样品；
  \item 不同间隔层厚度器件的系统对比；
  \item 不同调制频率与光强条件下的响应对比。
\end{itemize}

% ==================================================
\section{预期结果与创新点}

\subsection{预期结果}

\begin{itemize}
  \item 观测到与外加时间调制同步的光学强度变化；
  \item 在频域中出现清晰的调制频率峰及其谐波结构；
  \item 得到体系的调制带宽与时间常数；
  \item 揭示石墨烯–金纳米结构耦合对调制效率的影响规律。
\end{itemize}

\subsection{创新点}

\begin{itemize}
  \item 以时间调制为研究核心，而非静态 LSPR 传感；
  \item 将石墨烯作为主动调制单元引入等离激元体系；
  \item 从时间域与频域双重角度定量研究等离激元动态行为。
\end{itemize}

% ==================================================
\section{项目进度安排（示例）}

\begin{itemize}
  \item 阶段一：器件制备与静态光谱表征；
  \item 阶段二：时间调制实验与时间序列测量；
  \item 阶段三：频域分析与系统动力学建模；
  \item 阶段四：对照实验与机制归因总结。
\end{itemize}

\end{document}

